\documentclass{article}
\usepackage{amsmath}
\usepackage{amssymb}

\title{How to Write Mathematical Equations in LaTeX}
\author{}
\date{}

\begin{document}

\maketitle

\section{Inline Equations}
You can write equations inline like $E = mc^2$ or $a^2 + b^2 = c^2$ within your text.

\section{Display Equations}

\subsection{Basic Equation}
\begin{equation}
    f(x) = ax^2 + bx + c
    \label{eq:quadratic}
\end{equation}

\subsection{Quadratic Formula}
\begin{equation}
    x = \frac{-b \pm \sqrt{b^2 - 4ac}}{2a}
    \label{eq:quadratic_formula}
\end{equation}

\section{Aligned Equations}

\begin{align}
    (x+y)^2 &= x^2 + 2xy + y^2 \\
    (x-y)^2 &= x^2 - 2xy + y^2 \\
    x^2 - y^2 &= (x+y)(x-y)
\end{align}

\section{Matrices}

\subsection{Matrix Example}
\begin{equation}
    \begin{pmatrix}
        1 & 2 & 3 \\
        4 & 5 & 6 \\
        7 & 8 & 9
    \end{pmatrix}
\end{equation}

\section{Calculus}

\subsection{Derivatives}
\begin{equation}
    \frac{d}{dx}(x^n) = nx^{n-1}
\end{equation}

\subsection{Integrals}
\begin{equation}
    \int_a^b f(x) \, dx = F(b) - F(a)
\end{equation}

\subsection{Limits}
\begin{equation}
    \lim_{x \to \infty} \frac{1}{x} = 0
\end{equation}

\section{Complex Equations}

\begin{align}
    \sin^2(\theta) + \cos^2(\theta) &= 1 \\
    e^{i\pi} + 1 &= 0 \\
    \frac{\partial f}{\partial x} &= \lim_{h \to 0} \frac{f(x+h) - f(x)}{h}
\end{align}

\end{document}